% Section content template for individual Educative lessons
% This file contains the actual content for a specific section
% Generated from Educative API section content

% Section: Summary: Designing a Library Management System
% Section ID: 4646102778118144
% Section Slug: summary-designing-a-library-management-system
% Generated: 2025-12-11T15:13:38.818576

Now that you've completed the library management system case study, let's take a moment to reflect on and consolidate what we've learned. We 'll revisit the key system requirements, identify the core classes along with their responsibilities and relationships, and highlight the major design principles applied. We
'll also examine how objects interact within the system and walk through the overall workflow to understand how the components come together to achieve the desired functionality.

\subsection{Key requirements}\label{y6Ec3wVEecpm2HabgXzoK}
\\
This section outlines the primary functional requirements that define the scope and capabilities of the library management system.

\begin{enumerate}
\item
\protect\phantomsection\label{E0e4Mv41EXcNSpBeE1VU5}
Store and log all information about books, members, and transactions (borrowing, returns, reservations, renewals).
\item
\protect\phantomsection\label{2FnlvLQLyZUeJ55hdwNen}
Ensure every book has a unique ID and detailed information, including its physical location.
\item
\protect\phantomsection\label{qKTF5Ua5VyPJ6Uk3GGml0}
Include each book's ISBN, title, author, subject, and publication date.
\item
\protect\phantomsection\label{C4EsN9UHeox9osbiyhEHp}
Support each book title's multiple physical copies (book items).
\item
\protect\phantomsection\label{z6bz5WVmJlTcYEC7a_YZR}
Define two user roles: librarian and member, each with a unique library card.
\item
\protect\phantomsection\label{XsPqfVc9xmpdPGaSpIfcD}
Require every user to have a library card with a unique number.
\item
\protect\phantomsection\label{BsTG1re3UsRfRmbF-NLch}
Limit members to borrowing a maximum of 10 books at one time.
\item
\protect\phantomsection\label{hGjCag1fvoY9gC7k51-_b}
Set the maximum borrowing period for a book to 15 days.
\item
\protect\phantomsection\label{yk4tBno3dlZGtDR4sP2Iv}
Allow a book item to be reserved by only one member at a time.
\item
\protect\phantomsection\label{mRtzddxsrS84b1jXu2QZ5}
Record which user issued, reserved, or renewed a book item and the transaction date.
\item
\protect\phantomsection\label{s2KAp1nIjGcKeCe1mKRVr}
Enable members to renew borrowed books based on the library policy.
\item
\protect\phantomsection\label{wk-oEWGxL7brnWgHjl0AI}
Send notifications for overdue books and available reservations.
\item
\protect\phantomsection\label{B0jhsmGMVSAZ0sfBVXhKS}
Allow members to reserve a book if it is currently unavailable.
\item
\protect\phantomsection\label{OS3Whpbh1hTd_819t7G0d}
Enable users to search for books by title, author, subject, or publication date.
\end{enumerate}

\subsection{System actors}\label{S5RIL3SHqAlKmCkD9pNrV}
\\
This section identifies the primary and secondary actors who interact with the system.

\begin{itemize}
\item
\protect\phantomsection\label{bTz_RRhzEs3skV70Qahm8}
\textbf{Member:} A library user who can search the catalog, borrow, return, renew, and reserve books. They can also manage their account and pay fines.
\item
\protect\phantomsection\label{SgibUDA_5KPm-i9_T9d7V}
\textbf{Librarian:} An administrator responsible for managing the library's collection, including adding and removing books and book items. They also manage member accounts, issue books, and handle fines.
\end{itemize}

\subsection{Important classes and relationships}\label{zGgQngxslL8UdXx-j_Jy6}
\\
This section details the core classes that form the system's backbone, describing their purpose and how they relate to one another.

\begin{center}
\begin{tabular}{|p{0.18\linewidth}|p{0.42\linewidth}|p{0.35\linewidth}|}
\hline
\textbf{Class} & \textbf{Description} & \textbf{Important Relationships} \\
\hline
User & An abstract class representing a person interacting with the system. & Parent to Librarian and Member. \\
\hline
Librarian & A user with administrative privileges to manage books and members. & Inherits from User. Has an association with BookItem. \\
\hline
Member & A user who can borrow, reserve, and return books. & Inherits from User. \\
\hline
Book & Represents the conceptual information of a book (title, author, ISBN). & Has an association with Author. Aggregated in Catalog. \\
\hline
BookItem & Represents a specific physical copy of a book that can be loaned out. & Composition with Book. Has an association with Rack. \\
\hline
BookLending & Manages the checkout process, including due dates and return dates. & Has an association with User and BookItem. \\
\hline
BookReservation & Manages reservations for unavailable book items. & Has an association with User and BookItem. \\
\hline
Catalog & Implements the search functionality, allowing users to find books. & Implements the Search interface. Aggregates Book. \\
\hline
Notification & An abstract class for sending notifications to members. & Parent to EmailNotification and PostalNotification. \\
\hline
LibraryCard & Contains the user’s library card information. & Composition with User. \\
\hline
\end{tabular}
\end{center}


\subsection{Design highlights}\label{x6Yh0VcIE67hLz3CXYm8S}
\\
This section covers the overall design strategy, the core software design principles, and the key design patterns used to structure the system.

\subsubsection{Design approach}\label{tPPo_FH_5ZkphXcJ5ELjo}
\\
The system was designed using a bottom-up approach. This involved identifying core entities like \texttt{Book}, \texttt{Librarian}, and \texttt{Member}, modeling their workflows, and finally integrating these components into a complete, scalable system that adheres to SOLID principles.

\subsubsection{Design principles}\label{n-A_qf04hsHmycuNbRwXX}
\\
The design of the library management system incorporates the SOLID principles to ensure it is robust, maintainable, and scalable.

\begin{itemize}
\item
\protect\phantomsection\label{3bSejmhK9-OAwVhKGtfL_}
\textbf{Single Responsibility principle (SRP):} Each class has a distinct responsibility. For example, the \texttt{BookLending} class handles only the logic for borrowing books, while the \texttt{BookReservation} class manages only reservations.
\item
\protect\phantomsection\label{jlgwaav1J4iD4nMvlrp2N}
\textbf{Open/Closed principle (OCP):} The system is open for extension but closed for modification. For instance, new types of notifications (like SMS) can be added by creating a new subclass of \texttt{Notification} without altering existing code.
\item
\protect\phantomsection\label{RaxII7gDgjvuHpX8KjRPU}
\textbf{Liskov Substitution principle (LSP):} Subtypes can be substituted for their base types. A \texttt{Librarian} or \texttt{Member} object can be used wherever a \texttt{User} object is expected, as both inherit from \texttt{User}.
\item
\protect\phantomsection\label{yqfNoMYpf7iAiIxRN8dsL}
\textbf{Interface Segregation principle (ISP):} The \texttt{Search} interface is a good example, providing a specific contract for search functionality that the \texttt{Catalog} class implements, without forcing it to depend on unrelated methods.
\item
\protect\phantomsection\label{d9SZsSs-lXMkFXb1AwfHe}
\textbf{Dependency Inversion principle (DIP):} The system depends on abstractions, not concretions. The \texttt{Catalog} class depends on the \texttt{Search} interface, allowing different search implementations to be used interchangeably.
\end{itemize}

\subsubsection{Design patterns}\label{Bh0NVN1NdtnnrwFLN48kY}
\\
The system leverages several design patterns to solve common problems and improve the overall architecture.

\begin{itemize}
\item
\protect\phantomsection\label{CTWxXfBhRi4WxfxT-gSmB}
\textbf{Factory pattern:} Used to create domain objects like \texttt{Book} and \texttt{Member}, encapsulating the creation logic and ensuring objects are instantiated correctly.
\item
\protect\phantomsection\label{0vq0eUzY1YjYw8szK-Tvc}
\textbf{Delegation pattern:} Responsibilities are distributed among classes. A \texttt{Librarian} delegates the task of managing a book's status to the \texttt{BookItem} class.
\item
\protect\phantomsection\label{jRJmrpDZKI3-4_W5y_aoh}
\textbf{Observer pattern:} This pattern is used for notifications. Members can ``observe'' a \texttt{BookItem}; when its status changes (e.g., becomes available), they are automatically notified.
\end{itemize}

\subsection{Object interactions}\label{Kyl76gKTQg4hbVFJZTNSe}
\\
This section describes the system's dynamic behavior by explaining how key objects collaborate to perform tasks, based on the sequence diagram for issuing a book.

\begin{itemize}
\item
\protect\phantomsection\label{MsGSDzvFKEvH_UxUZwLTl}
\textbf{Book issuance request:} A \texttt{Member} initiates the process by requesting to borrow a \texttt{Book} from a \texttt{Librarian}.
\item
\protect\phantomsection\label{why5JovdadIAni9sBNEQH}
\textbf{Book status check:} If the member's quota is not exceeded, the \texttt{Librarian} checks the status of the \texttt{BookItem}.
\item
\protect\phantomsection\label{KuOtjTsGE5QSFr1UkVK-w}
\textbf{Completing the checkout:} If the \texttt{BookItem} is available, the \texttt{Librarian} issues the book to the \texttt{Member}. If the book is already reserved, the request is canceled.
\end{itemize}

\subsection{System workflow}\label{c8FhzkSfTii7jGoTZkcnH}
\\
This section outlines the primary end-to-end scenarios in the system, reflecting the logic from the activity diagrams.

\begin{itemize}
\item
\protect\phantomsection\label{jJkdUBWxySnjaEeGU6j9a}
\textbf{Checking out a book:} A member starts by providing their member ID and the book ID. The system verifies that the member's borrowing quota is not exceeded and that the book is not already reserved. If both checks pass, the book's status is updated to ``issued,'' and the member's count of borrowed books is incremented.
\item
\protect\phantomsection\label{ZncaIisDEs3uWKKm5pZV8}
\textbf{Returning a book:} A member provides the book ID to initiate a return. The system checks if the due date has passed. If it has, a fine is calculated and collected. The member's count of borrowed books is then decremented. Finally, the system checks if another member has reserved the book. If so, its status is changed to ``Reserved'' and the reserving member is notified; otherwise, it is set to ``Available.''
\end{itemize}

% End of section content